\chapter{Аудит кандидатов коэффициента гиперзарядового gap
\texorpdfstring{$R_2$}{R2}}
\label{ch:v10-h15-r2-hypercharge-gap-coefficient-audit}

Предыдущий гейт нашёл непустое окно знаков, но не выбрал его точку. Для
общего квадратичного полинома по $Q=6Y$
\begin{equation}
 H(a,b)=aI+bQ^2
\end{equation}
массы спектральных классов $Q^2=49,9,1$ равны
\begin{equation}
 m_{R_2}^2=a+49b,\qquad
 m_9^2=a+9b,\qquad
 m_1^2=a+b.
\end{equation}
Строгое выделение единственной тахионной $R_2$-пары эквивалентно
\begin{equation}
 \boxed{b<0,\qquad a+49b<0,\qquad a+9b>0.}
\end{equation}
После нормировки $b=-1$ получаем $(a,b)=(49-r,-1)$ и открытый интервал
$0<r<40$.

Этот конус не выбирает отношение. Две точные неколлинеарные нормированные
точки $(39,-1)$ и $(19,-1)$ соответствуют $r=10$ и $30$ и дают массы
\begin{equation}
 (-10,30,38),\qquad (-30,10,18).
\end{equation}
Обе проходят один и тот же sign gate. Напротив, сам алгебраический gap
$G_Y=49I-Q^2$ даёт $(0,40,48)$ и лежит на границе: $R_2$ остаётся
безмассовым, а не тахионным. Вторая граница $(9,-1)$ даёт $(-40,0,8)$.

Проверены одиннадцать механизмов по шести критериям: разрешение по $Q^2$,
внутренняя точка sign-конуса, положительная квартика, типизированный фон
нарушения до SM, унаследованный носитель и фиксация коэффициента. Матрица
имеет ранг $6$, но проходов нет: $0/11$.

Обычные spectral moments дают общий квадратично-квартичный масштаб, но не
$Q^2$-расщепление. Положительный Casimir или D-term поднимает $R_2$ сильнее
остальных; отрицательный Casimir делает все классы неустойчивыми.
Универсальный Coleman--Weinberg portal не меняет разности масс, а
dimensional transmutation способна задать масштаб, но не направление в
двумерном пространстве $(a,b)$.

Лучший кандидат с оценкой $5/6$ --- единый connected spectral trace на
типизированном фоне нарушения Pati--Salam до SM. Он мог бы одновременно
дать $Q^2$-чувствительность, внутренний знак, квартику и фиксированное
отношение. Провален ровно inherited-carrier criterion: текущая геометрия не
содержит общего динамического фона для $T^3_R$ и $B-L$, реализующего
$Y=T^3_R+(B-L)/2$ внутри того же родителя.

\begin{theorem}[Неединственность коэффициента gap]
Условия правильной инерции задают открытый двумерный конус коэффициентов.
Даже после удаления общей положительной нормировки остаётся интервал
$0<\mu^2/\kappa<40$. Поэтому спектр зарядов и требование знака не могут
вывести единственное отношение $\mu^2/\kappa$.
\end{theorem}

\begin{nogocriterion}[Target-loaded interior]
Добавление вручную любой точки $(49-r,-1)$ с $0<r<40$ доказывает лишь
существование подходящего потенциала. Физический parent требует, чтобы фон
нарушения, mixed invariant, квартика и их коэффициент возникли из одного
унаследованного оператора до проверки целевой сигнатуры.
\end{nogocriterion}

\begin{protocol}\begin{enumerate}
\item размерность коэффициентного пространства: \textbf{$2$};
\item sign-конус: \textbf{$b<0$, $a+49b<0$, $a+9b>0$};
\item точные внутренние отношения: \textbf{$10$ и $30$};
\item алгебраический $G_Y$: \textbf{граница $(0,40,48)$};
\item ранг матрицы кандидатов: \textbf{$6$};
\item прошедших кандидатов: \textbf{$0/11$};
\item лучший кандидат: \textbf{connected breaking-background trace, $5/6$};
\item inherited coefficient-map rank: \textbf{$0$};
\item physical origin: \textbf{$1/3$};
\item следующий гейт: \textbf{общий носитель гиперзарядового breaking-background}.
\end{enumerate}\end{protocol}

Машинный сертификат сохранён в
\path{s2t_v10_particle_wrinkle_dislocation_callias_equal_charge_twist_m4_h15_r2_hypercharge_gap_coefficient_candidate_audit_gate_results.json}.